\chapter{Blockchains as Adversarial State Machines}

\section{Purpose}

This chapter establishes the central mental model for the book:

\begin{quote}
A blockchain is a replicated state machine whose inputs are chosen in an
adversarial environment.
\end{quote}

That sentence is more useful than most slogans about decentralization,
immutability, or trustlessness. It gives us a way to reason about
security without getting trapped inside the vocabulary of any single
chain. Bitcoin, Ethereum, Solana, Move-based systems, appchains,
bridges, rollups, and smart contracts all differ in execution model and
consensus mechanism. But they all have to answer the same
first-principles questions:

\begin{itemize}
\tightlist
\item
  What is the state?
\item
  What inputs can change it?
\item
  Which transitions are valid?
\item
  Who gets to submit, order, censor, or finalize those inputs?
\item
  Which properties must remain true across every reachable state?
\item
  What assumptions make replication, progress, and agreement possible?
\end{itemize}

This chapter is deliberately not a tour of blockchain buzzwords. We will
not yet study signatures in depth, consensus protocols, token standards,
reentrancy, MEV, or bridges. Those topics come later. Here we build the
frame that makes those topics teachable.

If you learn one habit from this chapter, make it this:

\begin{quote}
Never review a blockchain system as a collection of features. Review it
as a set of possible state transitions under adversarial input and
ordering.
\end{quote}

\section{Learning Objectives}

By the end of this chapter, you should be able to:

\begin{itemize}
\tightlist
\item
  Explain a blockchain as a state machine rather than as a database or a
  payment network.
\item
  Distinguish state, transaction, block, transition function, validity
  rule, and invariant.
\item
  Explain why deterministic execution is necessary for replicated
  validation.
\item
  Identify where adversarial inputs enter a blockchain system.
\item
  Separate safety failures from liveness failures.
\item
  Explain why many serious exploits are not ``invalid transactions'' but
  valid transitions that violate protocol intent.
\end{itemize}

\section{The Core Model}

A state machine has four basic ingredients:

\begin{itemize}
\tightlist
\item
  A set of possible states.
\item
  A set of possible inputs.
\item
  A transition function that maps a current state and an input to a next
  state.
\item
  Rules for rejecting inputs or transitions that are not allowed.
\end{itemize}

In abstract form:

\begin{lstlisting}
next_state = apply(current_state, input)
\end{lstlisting}

For blockchains, the important inputs are usually transactions,
messages, blocks, votes, proofs, signatures, and governance actions. A
block is a batch of ordered inputs. A chain is a history of such
batches. A node validates the chain by starting from an initial state
and applying the accepted inputs in the accepted order.

Ethereum's formal specification explicitly frames Ethereum as a
transaction-based state machine: a genesis state is transformed into a
current state by executing transactions.\footnote{Gavin Wood,
  ``Ethereum: A Secure Decentralised Generalised Transaction Ledger'',
  https://ethereum.github.io/yellowpaper/paper.pdf.} Ethereum.org's
developer documentation makes the same operational point in simpler
terms: transactions are signed instructions from accounts that update
network state, and blocks batch transactions so participants can
synchronize around a common history.\footnote{ethereum.org,
  ``Transactions'', https://ethereum.org/developers/docs/transactions/.}
\footnote{ethereum.org, ``Blocks'',
  https://ethereum.org/developers/docs/blocks/.}

Bitcoin is often described as a ledger rather than a state machine, but
the same model applies. The state is the set of unspent transaction
outputs, along with consensus-relevant chain data. A transaction
consumes existing outputs and creates new outputs. A block proposes an
ordered set of such transactions. Nodes accept a block only if its
transactions satisfy the validity rules, including the rule that the
same output is not spent twice.\footnote{Satoshi Nakamoto, ``Bitcoin: A
  Peer-to-Peer Electronic Cash System'',
  https://bitcoin.org/bitcoin.pdf.}

This generalizes:

\begin{lstlisting}
genesis_state
  --tx_1--> state_1
  --tx_2--> state_2
  --tx_3--> state_3
  ...
\end{lstlisting}

For review work, draw the machine before naming bug classes:

\begin{lstlisting}
flowchart LR
    S[Current state] --> I[Input]
    I --> V{Validity check}
    V -- reject --> R[No state change]
    V -- accept --> T[Transition function]
    T --> N[Next state]
    N --> P{Invariant check}
    P -- holds --> A[Accept transition]
    P -- broken --> B[Unsafe reachable state]
\end{lstlisting}

The same idea can be written as minimal review pseudocode:

\begin{lstlisting}
apply(state, input):
    check syntax and encoding
    check authority
    check validity in current state
    compute next_state
    check required invariants
    return next_state or reject
\end{lstlisting}

For blocks:

\begin{lstlisting}
state_after_block = apply_block(parent_state, [tx_0, tx_1, tx_2, ...])
\end{lstlisting}

And for a chain:

\begin{lstlisting}
current_state = replay(genesis_state, canonical_history)
\end{lstlisting}

The security question is not merely whether a transaction is formatted
correctly. The question is whether the system's rules ensure that every
reachable state is acceptable under the intended security model.

That distinction will carry the whole book.

\section{State, State Transitions, and Validity Rules}

\subsection{State Is Everything the Rules Need to Remember}

State is the information that matters for future validity and behavior.

In a blockchain system, state may include:

\begin{itemize}
\tightlist
\item
  Account balances.
\item
  UTXO sets.
\item
  Contract storage.
\item
  Ownership records.
\item
  Nonces and sequence numbers.
\item
  Validator sets and stakes.
\item
  Governance configuration.
\item
  Pending withdrawals.
\item
  Oracle prices cached on-chain.
\item
  Token supplies.
\item
  Vault share balances.
\item
  Bridge message statuses.
\item
  Rollup output roots.
\item
  Protocol parameters such as fee rates, limits, and emergency flags.
\end{itemize}

Do not reduce state to ``balances.'' Balances are only one kind of
state. A paused flag is state. An upgrade admin is state. A stored
oracle timestamp is state. A consumed bridge message nonce is state. A
validator's slashing status is state. A record that a withdrawal has
already been claimed is state.

For security analysis, state is whatever the attacker wants to change,
depend on, corrupt, desynchronize, or freeze.

\subsection{Transitions Are the Only Way State Changes}

A state transition is a change from one state to another. In blockchain
systems, transitions are usually caused by accepted inputs:

\begin{itemize}
\tightlist
\item
  A user signs and submits a transaction.
\item
  A contract calls another contract.
\item
  A validator proposes a block.
\item
  A governance proposal executes.
\item
  A bridge relayer submits a message.
\item
  A sequencer posts a rollup batch.
\item
  A withdrawal proof is verified.
\item
  An oracle update is accepted.
\item
  A protocol upgrade changes the transition rules themselves.
\end{itemize}

The transition is not the transaction object. The transaction is the
input. The transition is the state change produced by applying the input
under the rules.

Consider a tiny token ledger:

{\def\LTcaptype{none} % do not increment counter
\begin{longtable}[]{@{}
  >{\raggedright\arraybackslash}p{(\linewidth - 2\tabcolsep) * \real{0.5000}}
  >{\raggedright\arraybackslash}p{(\linewidth - 2\tabcolsep) * \real{0.5000}}@{}}
\toprule\noalign{}
\begin{minipage}[b]{\linewidth}\raggedright
Element
\end{minipage} & \begin{minipage}[b]{\linewidth}\raggedright
Token ledger model
\end{minipage} \\
\midrule\noalign{}
\endhead
\bottomrule\noalign{}
\endlastfoot
State & \passthrough{\lstinline!balances[address]!},
\passthrough{\lstinline!total\_supply!} \\
Input &
\passthrough{\lstinline!transfer(sender, recipient, amount, signature)!} \\
Authority & \passthrough{\lstinline!signature!} must authorize
\passthrough{\lstinline!sender!} \\
Validity & \passthrough{\lstinline!amount > 0!} and
\passthrough{\lstinline!balances[sender] >= amount!} \\
Transition & Decrease sender, increase recipient, leave
\passthrough{\lstinline!total\_supply!} unchanged \\
Invariant & Sum of balances remains equal to
\passthrough{\lstinline!total\_supply!} \\
\end{longtable}
}

The transfer transition table is:

{\def\LTcaptype{none} % do not increment counter
\begin{longtable}[]{@{}
  >{\raggedright\arraybackslash}p{(\linewidth - 4\tabcolsep) * \real{0.3333}}
  >{\raggedright\arraybackslash}p{(\linewidth - 4\tabcolsep) * \real{0.3333}}
  >{\raggedright\arraybackslash}p{(\linewidth - 4\tabcolsep) * \real{0.3333}}@{}}
\toprule\noalign{}
\begin{minipage}[b]{\linewidth}\raggedright
Step
\end{minipage} & \begin{minipage}[b]{\linewidth}\raggedright
Check or update
\end{minipage} & \begin{minipage}[b]{\linewidth}\raggedright
Security purpose
\end{minipage} \\
\midrule\noalign{}
\endhead
\bottomrule\noalign{}
\endlastfoot
1 & Decode \passthrough{\lstinline!sender!},
\passthrough{\lstinline!recipient!}, \passthrough{\lstinline!amount!},
\passthrough{\lstinline!signature!} & Reject malformed inputs before
state changes \\
2 & Verify the signature authorizes \passthrough{\lstinline!sender!} &
Prevent unauthorized movement \\
3 & Check \passthrough{\lstinline!amount > 0!} & Reject nonsensical or
ambiguous movement \\
4 & Check \passthrough{\lstinline!balances[sender] >= amount!} & Prevent
negative balances \\
5 & Subtract from \passthrough{\lstinline!sender!} and add to
\passthrough{\lstinline!recipient!} & Move value without changing total
supply \\
6 & Recheck conservation & Catch implementation or integration errors \\
\end{longtable}
}

The transition is acceptable only if the validity rules hold. If
\passthrough{\lstinline!balances[sender]!} is less than
\passthrough{\lstinline!amount!}, the input must not produce the
transfer. If the signature does not authorize the sender, the input must
not move funds. If the transition decreases the sender but forgets to
increase the recipient, the transition violates asset conservation even
if the transaction looked well formed.

Security starts with this discipline:

\begin{lstlisting}
For every accepted input, what state changes?
For every rejected input, why is it rejected?
For every state change, what must remain true?
\end{lstlisting}

\subsection{Validity Is Not the Same as Safety}

The word ``valid'' is dangerous because it is often used at different
layers.

A transaction may be:

\begin{itemize}
\tightlist
\item
  Syntactically valid: it is encoded correctly.
\item
  Authenticated: the signature or authorization data checks out.
\item
  Executable: the current state has enough resources for execution.
\item
  Protocol-valid: it satisfies the base chain's validity rules.
\item
  Application-valid: it satisfies a smart contract or program's explicit
  checks.
\item
  Intent-safe: it preserves the economic or security properties the
  protocol was supposed to protect.
\end{itemize}

The last category is where many expensive failures live.

A lending protocol may accept a transaction that is syntactically valid,
signed by the attacker, included in a valid block, and allowed by the
contract code. Yet the same transaction may still leave the protocol
insolvent because the oracle price was manipulable or the collateral
accounting rounded in the attacker's favor. The chain did not fail to
validate the transaction. The application failed to encode the right
validity constraints or invariants.

This is why ``the transaction was valid'' is not a defense of a
protocol. It is often the description of the exploit.

\subsection{Validity Rules Define Reachable States}

The set of reachable states is the set of states the system can enter by
starting from the initial state and applying accepted transitions.

Informally:

\begin{lstlisting}
genesis_state is reachable

if state is reachable
and input is valid in state
then apply(state, input) is reachable
\end{lstlisting}

A security property must hold over reachable states, not imagined
states. If a state cannot occur under any accepted transition, it may be
irrelevant. If a dangerous state can occur through any valid sequence of
inputs, it matters even if the sequence looks unusual.

Attackers specialize in unusual valid sequences.

They will not use your protocol according to the product story. They
will use it according to the transition rules.

\subsection{Invariants Are Properties That Must Survive Every
Transition}

An invariant is a property that should remain true across all reachable
states.

Examples:

\begin{itemize}
\tightlist
\item
  Total token balances equal total supply.
\item
  No account balance is negative.
\item
  A withdrawal cannot be claimed twice.
\item
  A user cannot spend funds without authorization.
\item
  A vault's assets are sufficient to redeem outstanding shares according
  to the redemption rules.
\item
  A bridge message cannot be executed on the destination chain unless it
  was committed on the source chain under the required finality
  assumptions.
\item
  A governance executor cannot bypass the configured delay unless
  emergency powers explicitly allow it.
\end{itemize}

The invariant mindset comes from state-machine reasoning. Lamport's
treatment of state machines emphasizes invariants as predicates that
hold in all reachable states, and proves them by showing they hold
initially and are preserved by transitions.\footnote{Leslie Lamport,
  ``Computation and State Machines'',
  https://lamport.azurewebsites.net/pubs/state-machine.pdf.} This is not
academic ornamentation. It is exactly the reasoning security reviewers
need.

For a blockchain protocol, the practical version is:

\begin{lstlisting}
If the invariant is true before this transaction,
is it still true after this transaction?
\end{lstlisting}

If the answer is no, you have either found a bug or discovered that your
invariant is not actually part of the intended design.

Both outcomes are valuable.

\subsection{State Includes the Rules That Can Change the Rules}

Many blockchain systems have upgrade mechanisms, governance modules,
feature flags, admin keys, pausing controls, validator-set updates, and
parameter changes. These are not external to the state machine. They are
transitions that alter future transition rules or future authority.

For example:

\begin{lstlisting}
State:
  implementation_address
  admin
  timelock_delay
  paused

Input:
  upgrade_to(new_implementation)

Validity:
  caller == admin
  timelock elapsed
  new implementation passes required checks

Transition:
  implementation_address = new_implementation
\end{lstlisting}

An upgrade does not merely change a value. It changes the code that
determines future state transitions. That makes upgrade authority one of
the most important security boundaries in any system that has it.

Do not treat governance and upgrades as operational details. They are
part of the state machine.

\section{Determinism, Replication, and Global Agreement}

\subsection{Replication Requires Determinism}

A blockchain is replicated. Many nodes maintain or verify the same
system. They may be run by different people, on different machines, in
different countries, using different client implementations.

Replication only works if honest validators can independently reach the
same result.

If two honest nodes start from the same state and apply the same valid
input sequence, they must compute the same next state. Otherwise the
network does not have one replicated state machine. It has a collection
of machines drifting apart.

This requirement can be stated simply:

\begin{lstlisting}
same initial state
+ same ordered inputs
+ same transition rules
= same resulting state
\end{lstlisting}

Determinism is why blockchain execution environments restrict or
carefully define sources of environmental data. A contract cannot ask an
arbitrary web server for the current weather during consensus execution.
A validator cannot use its local machine clock as an unbounded source of
truth. A program cannot depend on uncommitted local randomness and still
expect all validators to agree.

When a blockchain exposes ``environmental'' values such as block number,
timestamp, slot, block hash, base fee, recent blockhash, or chain ID,
those values are not private local facts. They are consensus-visible
inputs or protocol-defined values. They are part of the transition
context.

This has a security consequence:

\begin{quote}
If a value is available to on-chain execution, assume it is public,
constrained by protocol rules, and potentially influenced by actors who
participate in block production.
\end{quote}

That does not mean every such value is useless. It means the security
assumption must be explicit.

\subsection{Consensus Orders Inputs; Execution Applies Them}

For teaching purposes, separate two jobs:

\begin{itemize}
\tightlist
\item
  Consensus, fork choice, or block production determines which ordered
  history is canonical.
\item
  Execution applies the transition rules to that ordered history.
\end{itemize}

Different chains combine these jobs differently, but the distinction
matters.

Bitcoin's proof-of-work chain selection gives nodes a way to converge on
a history of valid blocks under a resource assumption. Ethereum's modern
proof-of-stake protocol uses validators, attestations, fork choice, and
finality checkpoints. BFT-style appchains may use voting rounds and
finality certificates. Rollups may use sequencers, L1 contracts, fraud
proofs, validity proofs, or data availability assumptions.

Those mechanisms differ. The state-machine question remains:

\begin{lstlisting}
Which ordered inputs are accepted as canonical?
What state results from applying them?
When can users treat that result as settled?
\end{lstlisting}

Chapter 8 will study consensus and finality in detail. For now, do not
confuse ``valid transition'' with ``canonical transition.'' A
transaction can be locally valid and still not land in the canonical
chain. A block can contain valid transactions and still lose a
fork-choice race. A withdrawal can be valid on an L2 but not yet settled
on L1.

Security depends on both execution validity and history selection.

\subsection{Blocks Are Commitment Devices}

A block is not merely a container. It commits to an ordered set of data
and to a relationship with prior history.

Ethereum documentation describes blocks as batches of transactions
linked by hashes of previous blocks, so that changing prior data affects
later commitments.\footnote{ethereum.org, ``Blocks'',
  https://ethereum.org/developers/docs/blocks/.} The Ethereum Yellow
Paper makes a more formal point: blocks record transactions, reference
the prior block, and include an identifier of the resulting state rather
than storing the entire state itself.\footnote{Gavin Wood, ``Ethereum: A
  Secure Decentralised Generalised Transaction Ledger'',
  https://ethereum.github.io/yellowpaper/paper.pdf.}

This matters for security because nodes do not need to trust a block
producer's claim about the final state. They can execute the block and
verify that the resulting state commitment is correct.

At a high level:

\begin{lstlisting}
proposed block:
  parent reference
  ordered transactions
  state commitment
  other consensus data

validator checks:
  parent is acceptable
  transactions are valid in order
  execution result matches commitment
  consensus rules are satisfied
\end{lstlisting}

If the block producer lies about the state result, honest validating
nodes should reject the block. If clients disagree about the transition
rules, a chain split can occur. If the commitment scheme has a flaw,
proofs and light-client assumptions can fail. Those are different
failure modes, but all are visible in the state-machine frame.

\subsection{Global Agreement Is Not Instant Agreement}

``The blockchain state'' sounds singular. In reality, nodes observe the
network through time, latency, forks, local mempools, peer choices, RPC
providers, and client software. Agreement is a property achieved under
assumptions, not magic.

At a given moment:

\begin{itemize}
\tightlist
\item
  Two nodes may have seen different pending transactions.
\item
  Two nodes may have different mempool contents.
\item
  Two nodes may temporarily prefer different chain heads.
\item
  A block that appears confirmed may later be reorged in
  probabilistic-finality systems.
\item
  A rollup transaction may be accepted by a sequencer before it is
  settled on the base layer.
\item
  A user interface may show state from an indexer that is behind, wrong,
  or pointed at a minority fork.
\end{itemize}

Security reviews must be precise about which level of agreement is being
assumed:

\begin{itemize}
\tightlist
\item
  Local execution result.
\item
  Mempool observation.
\item
  Included in a block.
\item
  Confirmed after some depth.
\item
  Economically finalized.
\item
  BFT-finalized.
\item
  Settled on another domain.
\item
  Reflected correctly by an RPC, indexer, or frontend.
\end{itemize}

Many application bugs come from treating one level as if it were
another.

Example: a bridge that releases funds on chain B after seeing an event
on chain A must define how final that event must be. Seeing a
transaction in a block is not the same as accepting that the event can
never be removed from the canonical history. If the bridge treats
tentative state as final state, the bridge has imported the source
chain's reorg risk into its own asset accounting.

\subsection{Determinism Does Not Eliminate Implementation Risk}

A protocol can be deterministic on paper and still fail in
implementation.

Possible causes include:

\begin{itemize}
\tightlist
\item
  Client bugs.
\item
  Parser differences.
\item
  Integer overflow or underflow in a client implementation.
\item
  Inconsistent handling of malformed data.
\item
  Different interpretations of serialization rules.
\item
  Non-deterministic iteration over maps or sets.
\item
  Platform-specific behavior.
\item
  Incorrect gas or resource accounting.
\item
  Database corruption or state-pruning errors.
\end{itemize}

This is why blockchain security is not only smart contract security. The
replicated state machine depends on correct clients, correct
specifications, compatible implementations, and robust handling of
adversarial data.

If a malformed block causes one client to accept and another to reject,
the failure is not at the application layer. The replicated machine has
split on the meaning of its transition rules.

\section{Permissionless Inputs and Adversarial Execution}

\subsection{Permissionless Means Attackers Are Normal Users}

In many traditional systems, a backend receives requests from
authenticated users through controlled interfaces. Attackers may still
send malicious requests, but the system often has additional
chokepoints: web servers, rate limits, IP-based controls, account
review, fraud teams, manual rollback, or privileged database access.

Public blockchains change the default.

If a function, instruction, transaction type, or message path is public,
an attacker can usually call it. If a contract can be composed with, an
attacker can compose with it. If a state can be read, an attacker can
simulate against it. If a transaction can be frontrun, copied,
sandwiched, delayed, or backrun, someone may try.

The attacker is not outside the system. The attacker is a participant
using the system's own rules.

This is the permissionless input model:

\begin{lstlisting}
Any allowed actor may submit any syntactically valid input,
at any allowed time,
with any allowed parameters,
in any allowed sequence,
to any public entry point,
as long as they can pay the required costs.
\end{lstlisting}

That is the baseline. Then you add the chain-specific details: mempools,
fees, account models, signer semantics, compute limits, validator
behavior, sequencer rules, relayers, governance, and off-chain
infrastructure.

\subsection{Attackers Choose Edge Cases, Not Happy Paths}

A user story says:

\begin{lstlisting}
Alice deposits collateral, borrows conservatively, repays later, and withdraws.
\end{lstlisting}

An attacker asks:

\begin{lstlisting}
Can I deposit and withdraw in the same transaction?
Can I manipulate the price before borrowing?
Can I enter with zero shares?
Can I donate assets directly to skew accounting?
Can I call through a contract that reenters?
Can I create many accounts?
Can I force rounding to favor me?
Can I make another user's operation fail?
Can I trigger a state where exits are impossible?
Can I exploit a stale cache?
Can I submit the same message on another chain?
Can I reorder operations so the invariant fails?
\end{lstlisting}

The difference is not attitude. It is method.

Security review is adversarial state exploration. It asks what the
transition system permits, not what the designers expected honest users
to do.

\subsection{Public State Changes the Defender's Burden}

On-chain systems often expose state, code, events, pending transactions,
liquidity, prices, governance proposals, and privileged roles. Attackers
can inspect them cheaply. They can simulate transactions before sending
them. They can fork mainnet locally. They can watch pending
transactions. They can build bots that react faster than humans.

Public state has benefits: auditability, verifiability, composability,
and independent monitoring. But it also changes the defensive burden.

Do not rely on:

\begin{itemize}
\tightlist
\item
  Hidden contract logic.
\item
  Attackers not knowing the current state.
\item
  Attackers not finding a profitable ordering.
\item
  Attackers not composing your protocol with another.
\item
  Attackers not noticing a governance proposal.
\item
  Attackers not simulating the outcome.
\item
  Attackers behaving like UI users.
\end{itemize}

In public systems, obscurity is usually temporary. Security must live in
the transition rules, economic assumptions, and operational controls.

\subsection{The Adversary Controls More Than Transaction Parameters}

At minimum, an attacker controls the transactions they submit. Often
they can influence more:

\begin{itemize}
\tightlist
\item
  Timing: when to submit, delay, replace, or cancel.
\item
  Ordering: through priority fees, private order flow, builder
  relationships, or sequencer influence.
\item
  Atomicity: bundling multiple operations in one transaction.
\item
  Call graph: routing through contracts or programs that trigger
  callbacks.
\item
  Identity surface: creating many accounts or addresses.
\item
  Liquidity: moving prices temporarily.
\item
  Governance power: buying, borrowing, or bribing votes.
\item
  Infrastructure view: choosing which RPC, relayer, indexer, or frontend
  to interact with.
\item
  Cross-domain timing: exploiting finality gaps between chains or
  layers.
\end{itemize}

The exact powers differ by system. The habit is the same: make the
adversary's powers explicit.

Bad security reasoning says:

\begin{lstlisting}
No rational user would do that.
\end{lstlisting}

Better security reasoning says:

\begin{lstlisting}
Can any valid actor do that, and would doing it improve their position?
\end{lstlisting}

\subsection{Rejection Is a Security Boundary}

A blockchain system is defined as much by what it rejects as by what it
accepts.

Rejection can happen at several layers:

\begin{itemize}
\tightlist
\item
  A peer rejects malformed network data.
\item
  A node rejects an invalid transaction.
\item
  A block builder refuses to include a transaction.
\item
  A validator rejects an invalid block.
\item
  A contract reverts.
\item
  A program returns an error.
\item
  A bridge rejects a proof.
\item
  A governance executor rejects a payload that has not passed the
  timelock.
\end{itemize}

For reviewers, every rejection path deserves attention:

\begin{itemize}
\tightlist
\item
  Is the rejection condition complete?
\item
  Is it checked before any irreversible side effect?
\item
  Does failure revert all state changes or only some?
\item
  Can the rejection be triggered to grief others?
\item
  Is rejection deterministic across validators?
\item
  Does the caller learn enough to adapt and exploit?
\item
  Does the system fail closed or fail open?
\end{itemize}

In many smart contract exploits, the missing line is a missing rejection
condition:

\begin{lstlisting}
require(caller == owner)
require(amount <= balance)
require(message_not_already_used)
require(price_is_fresh)
require(account.owner == expected_program)
require(output_root_is_final)
\end{lstlisting}

But a rejection condition is not automatically correct just because it
exists. It may check the wrong authority, use stale state, validate the
wrong account, trust the wrong domain, or run after an external call.

\subsection{Adversarial Execution Includes Honest-Looking Transactions}

Not every harmful transaction is obviously malicious.

A transaction can be:

\begin{itemize}
\tightlist
\item
  Signed by the legitimate key.
\item
  Routed through a normal UI.
\item
  Included by an honest validator.
\item
  Valid under the protocol rules.
\item
  Economically profitable because of an unintended state interaction.
\end{itemize}

This is common in wallet drainers, approval phishing, governance
manipulation, oracle manipulation, liquidation races, and MEV
extraction. The security failure may not be ``someone hacked
consensus.'' It may be that the system allowed a user, bot, admin,
relayer, or contract to authorize a transition whose consequences were
not understood.

This is why the next chapter treats assets, ownership, authority, and
trust boundaries as first-class topics.

\section{Safety, Liveness, and Fault Assumptions}

\subsection{Safety: Bad States Must Not Happen}

Safety properties say that certain bad things must never happen. In
blockchain security, safety is about preventing unacceptable states or
conflicting histories.

Examples:

\begin{itemize}
\tightlist
\item
  No one can spend the same coin twice in the finalized history.
\item
  No account can transfer funds without authorization.
\item
  No accepted block can create coins outside the issuance rules.
\item
  No bridge message can be executed twice.
\item
  No finalized checkpoint conflicts with another finalized checkpoint
  under the protocol's assumptions.
\item
  No vault can issue claims that exceed redeemable assets under its
  accounting rules.
\item
  No governance action can execute before the required delay.
\end{itemize}

Safety failures are usually severe because the bad state has already
happened. A later transaction may repair accounting, compensate users,
or pause the system, but the safety violation itself is part of history
unless the chain or application has an explicit recovery mechanism.

In state-machine terms, safety asks:

\begin{lstlisting}
Can any valid sequence of transitions reach a forbidden state?
\end{lstlisting}

If yes, the system is unsafe with respect to that property.

\subsection{Liveness: Necessary Progress Must Remain Possible}

Liveness properties say that something required can eventually happen.

Examples:

\begin{itemize}
\tightlist
\item
  Valid transactions can eventually be included.
\item
  A user can eventually withdraw funds.
\item
  A rollup user can eventually force inclusion or exit.
\item
  A bridge message can eventually be relayed.
\item
  A liquidation system can eventually clear unhealthy debt.
\item
  A governance proposal can eventually execute after approval.
\item
  A paused protocol can eventually recover through a defined path.
\end{itemize}

Liveness failures can be just as destructive as safety failures. A
protocol that never loses accounting correctness but traps user funds
forever is not secure in any meaningful sense. A chain that preserves
safety by halting may be preferable to accepting conflicting histories,
but the halt is still a liveness failure with economic consequences.

In state-machine terms, liveness asks:

\begin{lstlisting}
From an allowed state, can the system still make required progress
under the stated assumptions?
\end{lstlisting}

If no, the system is live only on paper.

\subsection{Safety and Liveness Pull Against Each Other}

Distributed systems often face a tradeoff between preserving safety and
maintaining liveness under adverse network conditions.

If the network is partitioned, a system may choose to stop finalizing
rather than finalize conflicting histories. That preserves safety but
sacrifices liveness. A system that continues accepting conflicting
histories may preserve local availability but violate global safety.

Blockchain designs make different choices:

\begin{itemize}
\tightlist
\item
  Nakamoto-style systems may continue producing competing branches and
  resolve them probabilistically through fork choice.
\item
  BFT-style systems may stop finalizing if they cannot gather enough
  votes.
\item
  Rollups may allow users to exit through L1 mechanisms if the sequencer
  stops cooperating, but only if the escape hatch is actually
  implemented and usable.
\item
  Bridges may wait for deeper finality before releasing funds,
  sacrificing speed to reduce reorg risk.
\end{itemize}

The security reviewer should not ask ``Is it decentralized?'' as a vague
comfort phrase. Ask:

\begin{lstlisting}
Under which network and adversary assumptions is safety preserved?
Under which assumptions is progress guaranteed?
What happens outside those assumptions?
\end{lstlisting}

\subsection{Fault Assumptions Are Part of the Security Model}

No blockchain system is secure without assumptions. The assumptions may
be cryptographic, economic, network-level, implementation-level, or
human.

Common assumptions include:

\begin{itemize}
\tightlist
\item
  Hash functions are collision resistant for practical adversaries.
\item
  Signature schemes cannot be forged.
\item
  Honest nodes execute the protocol rules correctly.
\item
  A sufficient fraction of block production power follows the protocol.
\item
  The network eventually delivers messages.
\item
  Validators with enough stake are not colluding to violate safety.
\item
  At least one honest relayer exists for liveness.
\item
  A sequencer cannot permanently censor without users escaping through
  L1.
\item
  A multisig threshold is not compromised.
\item
  An oracle committee reports fresh and accurate prices within defined
  bounds.
\item
  A client bug does not affect a supermajority of economic weight.
\item
  Governance voters have enough time and information to react to
  malicious proposals.
\end{itemize}

These are not decorative. They determine the boundary of the security
claim.

Weak writing says:

\begin{lstlisting}
The protocol is secure because it is decentralized.
\end{lstlisting}

Serious writing says:

\begin{lstlisting}
The protocol preserves safety if fewer than one-third of voting power is Byzantine during the relevant finality window, clients agree on the transition rules, and the network eventually delivers messages. It preserves withdrawal liveness if at least one honest actor can submit the exit transaction and the base layer remains live.
\end{lstlisting}

The second statement can be attacked, tested, debated, and improved. The
first cannot.

\subsection{Byzantine, Rational, and Accidental Faults Are Different}

Faults are not all the same.

A crash fault occurs when a component stops. A Byzantine fault occurs
when a component behaves arbitrarily: lying, equivocating, censoring,
fabricating data, or behaving inconsistently toward different peers. A
rational adversary behaves strategically to maximize payoff. An
accidental fault may come from a bug, misconfiguration, operator error,
or misunderstood integration.

Blockchain systems often face all of these at once:

\begin{itemize}
\tightlist
\item
  A validator can go offline.
\item
  A validator can equivocate.
\item
  A block builder can reorder transactions for profit.
\item
  A user can sign a malicious approval by mistake.
\item
  A developer can deploy with the wrong address.
\item
  A client can parse a malformed block incorrectly.
\item
  A governance participant can vote for a payload they did not inspect.
\item
  A relayer can stop submitting messages because fees are uneconomic.
\end{itemize}

Do not collapse these into one generic ``attacker.'' Each fault type
requires different controls.

\begin{itemize}
\tightlist
\item
  Crash faults need redundancy and recovery.
\item
  Byzantine faults need validation, quorum assumptions, slashing, fraud
  proofs, or cryptographic proofs.
\item
  Rational faults need incentive analysis and payoff bounds.
\item
  Human faults need operational controls, interfaces, delays, and review
  processes.
\end{itemize}

\subsection{Time and Finality Are Security Parameters}

Many blockchain systems are secure only relative to time.

Examples:

\begin{itemize}
\tightlist
\item
  A Bitcoin payment becomes harder to reverse as more blocks are built
  after it, under the assumption that the attacker controls less hash
  power than the honest network.\footnote{Satoshi Nakamoto, ``Bitcoin: A
    Peer-to-Peer Electronic Cash System'',
    https://bitcoin.org/bitcoin.pdf.}
\item
  A proof-of-stake chain may justify or finalize checkpoints after a
  voting process.
\item
  A rollup withdrawal may require a challenge window.
\item
  An oracle price may be safe only if it is fresh.
\item
  A governance delay may be safe only if users can detect and respond
  before execution.
\item
  A liquidation may be safe only if keepers can act before bad debt
  accumulates.
\end{itemize}

Security review should therefore ask:

\begin{lstlisting}
When is this state considered final enough for the next action?
How long can an attacker delay progress?
How stale can this input become before it is unsafe?
How much time do defenders have to react?
\end{lstlisting}

Time is not an implementation detail. It is part of the threat model.

\section{Security Failures as Broken State Transitions}

\subsection{The Most Important Distinction}

There are two broad classes of state-machine failure:

\begin{enumerate}
\def\labelenumi{\arabic{enumi}.}
\tightlist
\item
  The system accepts a transition that should be invalid under its
  explicit rules.
\item
  The system accepts a transition that is valid under its explicit rules
  but violates the intended security properties.
\end{enumerate}

The first class is usually a protocol or implementation failure.
Examples include a signature verification bug, a consensus client bug, a
parser discrepancy, or an invalid block accepted by validators.

The second class is the center of application-layer blockchain security.
The chain does exactly what the code says. The code does not say what
the protocol actually needed.

Most smart contract exploits are in the second class.

That is why this book will repeatedly separate:

\begin{lstlisting}
rule validity
from
security intent
\end{lstlisting}

\subsection{A Taxonomy of Broken Transitions}

The state-machine model gives us a clean failure taxonomy.

{\def\LTcaptype{none} % do not increment counter
\begin{longtable}[]{@{}
  >{\raggedright\arraybackslash}p{(\linewidth - 4\tabcolsep) * \real{0.3333}}
  >{\raggedright\arraybackslash}p{(\linewidth - 4\tabcolsep) * \real{0.3333}}
  >{\raggedright\arraybackslash}p{(\linewidth - 4\tabcolsep) * \real{0.3333}}@{}}
\toprule\noalign{}
\begin{minipage}[b]{\linewidth}\raggedright
Failure type
\end{minipage} & \begin{minipage}[b]{\linewidth}\raggedright
State-machine description
\end{minipage} & \begin{minipage}[b]{\linewidth}\raggedright
Example
\end{minipage} \\
\midrule\noalign{}
\endhead
\bottomrule\noalign{}
\endlastfoot
Invalid transition accepted & The validator accepts a transition outside
the rules & Invalid block, forged signature, client consensus bug \\
Valid transition violates intent & The rules allow a reachable state
that breaks an invariant & Accounting bug, oracle manipulation, share
inflation \\
Unauthorized transition & State changes under the wrong authority &
Missing owner check, wrong signer, account substitution \\
Duplicate transition & A one-time transition executes more than once &
Replay, double claim, bridge message reused \\
Missing transition & A required state change cannot happen & Stuck
withdrawal, frozen queue, no recovery path \\
Wrong ordering & The same inputs in a different order produce harmful
state & Front-running, sandwiching, liquidation race \\
Divergent replicas & Honest nodes compute different states from the same
history & Client bug, nondeterminism, parser inconsistency \\
Unsafe external assumption & Transition relies on a fact not guaranteed
by the system & Stale oracle, unfinalized source-chain event,
compromised relayer \\
\end{longtable}
}

This taxonomy is not exhaustive, but it is a better starting point than
memorizing bug names. Bug names are symptoms. Broken transitions are
causes.

\subsection{Valid Exploits Are Still Exploits}

Suppose a vault has this simplified design:

\begin{lstlisting}
State:
  asset_balance
  total_shares
  shares[user]

deposit(amount):
  if total_shares == 0:
    minted = amount
  else:
    minted = amount * total_shares / asset_balance

  transfer amount from user to vault
  shares[user] += minted
  total_shares += minted

withdraw(shares_to_burn):
  amount = shares_to_burn * asset_balance / total_shares
  shares[user] -= shares_to_burn
  total_shares -= shares_to_burn
  transfer amount from vault to user
\end{lstlisting}

At first glance, the transition rules look reasonable. Deposits mint
shares. Withdrawals burn shares. The exchange rate comes from assets
divided by shares.

Now ask state-machine questions:

\begin{itemize}
\tightlist
\item
  Can \passthrough{\lstinline!asset\_balance!} change without
  \passthrough{\lstinline!total\_shares!} changing?
\item
  What happens when \passthrough{\lstinline!total\_shares!} is very
  small?
\item
  Which way does integer division round?
\item
  Can a deposit mint zero shares?
\item
  Does the first depositor get special power over the exchange rate?
\item
  Is direct donation to the vault part of the transition model?
\end{itemize}

If an attacker can become the first depositor with a tiny amount, then
donate assets directly to increase
\passthrough{\lstinline!asset\_balance!}, they may manipulate the share
price so later depositors receive too few shares because of rounding.
Every individual transaction may be valid. The exploit is a valid path
through the state machine that violates the intended fairness and
accounting properties.

The lesson is not ``all vaults are broken.'' The lesson is that security
properties must be encoded and tested against all reachable states,
including states produced by unusual but valid sequences.

\subsection{The Chain Protects Its Own Rules, Not Your Intent}

Base-layer validation enforces base-layer rules. It does not know your
business logic unless your business logic is encoded in executable
rules.

A chain can verify:

\begin{itemize}
\tightlist
\item
  This signature is valid.
\item
  This nonce is correct.
\item
  This transaction pays enough fee.
\item
  This block follows consensus rules.
\item
  This contract execution ended in a valid state according to the VM.
\end{itemize}

The chain generally cannot verify:

\begin{itemize}
\tightlist
\item
  This lending market's risk parameters are economically safe.
\item
  This oracle reflects a fair market price.
\item
  This governance proposal is not malicious.
\item
  This bridge message corresponds to a source event that your protocol
  should trust.
\item
  This user understood the approval they signed.
\item
  This vault share price is fair to later depositors.
\end{itemize}

Those properties must be designed, constrained, monitored, or explicitly
accepted as trust assumptions.

\subsection{Reorgs and Ordering Are State-Transition Problems}

A reorg is not only a chain-layer curiosity. It changes which state
transitions are part of canonical history.

Suppose a protocol observes:

\begin{lstlisting}
Block A includes deposit.
Protocol credits user on another domain.
Block A is reorged out.
Deposit no longer exists in canonical history.
Credit remains on other domain.
\end{lstlisting}

The failure is a mismatch between state machines. One domain treated a
transition as final before the other domain's finality assumptions
justified it.

Ordering has similar consequences. Consider two transactions:

\begin{lstlisting}
tx_1: oracle price update
tx_2: borrow against collateral
\end{lstlisting}

The resulting state may differ depending on whether
\passthrough{\lstinline!tx\_1!} happens before
\passthrough{\lstinline!tx\_2!} or after it. If attackers can influence
ordering, they can sometimes choose which state transition path the
system takes.

That is why mempools, block production, finality, bridges, and MEV
belong in a blockchain security curriculum. They are not separate from
smart contract security. They determine the input order and history that
the state machine accepts.

\subsection{Security Review Starts With Transition Questions}

When you encounter a new blockchain system, start with these questions:

\begin{enumerate}
\def\labelenumi{\arabic{enumi}.}
\tightlist
\item
  What is the state?
\item
  What are the entry points that can change it?
\item
  What are the validity rules for each entry point?
\item
  Who or what provides authority for each transition?
\item
  What external facts does the transition trust?
\item
  What ordering assumptions does the transition rely on?
\item
  What must remain true after every transition?
\item
  What required transitions must remain possible?
\item
  What happens if the input is delayed, duplicated, reordered, replayed,
  or removed from history?
\item
  What happens if an authorized actor behaves maliciously or
  irrationally?
\end{enumerate}

These questions are intentionally repetitive. Repetition is how you
avoid being fooled by surface complexity.

\section{Worked Example: A Minimal Escrow State Machine}

Consider a simple escrow:

\begin{lstlisting}
State:
  buyer
  seller
  arbiter
  amount
  funded: boolean
  released: boolean
  refunded: boolean

Transitions:
  fund()
  release()
  refund()
  arbitrate_release()
  arbitrate_refund()
\end{lstlisting}

The state diagram is:

\begin{lstlisting}
stateDiagram-v2
    [*] --> Unfunded
    Unfunded --> Funded: fund
    Funded --> Released: release / arbitrate_release
    Funded --> Refunded: refund / arbitrate_refund
    Released --> [*]
    Refunded --> [*]
\end{lstlisting}

A product description might say:

\begin{lstlisting}
The buyer funds the escrow. If delivery succeeds, the seller is paid.
If delivery fails, the buyer is refunded. The arbiter resolves disputes.
\end{lstlisting}

That is not yet a security model.

A state-machine review asks for explicit transition rules:

{\def\LTcaptype{none} % do not increment counter
\begin{longtable}[]{@{}
  >{\raggedright\arraybackslash}p{(\linewidth - 8\tabcolsep) * \real{0.2000}}
  >{\raggedright\arraybackslash}p{(\linewidth - 8\tabcolsep) * \real{0.2000}}
  >{\raggedright\arraybackslash}p{(\linewidth - 8\tabcolsep) * \real{0.2000}}
  >{\raggedright\arraybackslash}p{(\linewidth - 8\tabcolsep) * \real{0.2000}}
  >{\raggedright\arraybackslash}p{(\linewidth - 8\tabcolsep) * \real{0.2000}}@{}}
\toprule\noalign{}
\begin{minipage}[b]{\linewidth}\raggedright
Transition
\end{minipage} & \begin{minipage}[b]{\linewidth}\raggedright
Caller
\end{minipage} & \begin{minipage}[b]{\linewidth}\raggedright
Preconditions
\end{minipage} & \begin{minipage}[b]{\linewidth}\raggedright
State change
\end{minipage} & \begin{minipage}[b]{\linewidth}\raggedright
Value movement
\end{minipage} \\
\midrule\noalign{}
\endhead
\bottomrule\noalign{}
\endlastfoot
\passthrough{\lstinline!fund()!} & \passthrough{\lstinline!buyer!} &
\passthrough{\lstinline!funded == false!} and payment received &
\passthrough{\lstinline!funded = true!} & Buyer deposits
\passthrough{\lstinline!amount!} into escrow \\
\passthrough{\lstinline!release()!} & \passthrough{\lstinline!buyer!} &
\passthrough{\lstinline!funded == true!},
\passthrough{\lstinline!released == false!},
\passthrough{\lstinline!refunded == false!} &
\passthrough{\lstinline!released = true!} & Escrow pays seller \\
\passthrough{\lstinline!refund()!} & \passthrough{\lstinline!seller!} &
\passthrough{\lstinline!funded == true!},
\passthrough{\lstinline!released == false!},
\passthrough{\lstinline!refunded == false!} &
\passthrough{\lstinline!refunded = true!} & Escrow pays buyer \\
\passthrough{\lstinline!arbitrate\_release()!} &
\passthrough{\lstinline!arbiter!} &
\passthrough{\lstinline!funded == true!},
\passthrough{\lstinline!released == false!},
\passthrough{\lstinline!refunded == false!} &
\passthrough{\lstinline!released = true!} & Escrow pays seller \\
\passthrough{\lstinline!arbitrate\_refund()!} &
\passthrough{\lstinline!arbiter!} &
\passthrough{\lstinline!funded == true!},
\passthrough{\lstinline!released == false!},
\passthrough{\lstinline!refunded == false!} &
\passthrough{\lstinline!refunded = true!} & Escrow pays buyer \\
\end{longtable}
}

Now write invariants:

\begin{lstlisting}
released and refunded cannot both be true

if released == true:
  seller has been paid exactly once

if refunded == true:
  buyer has been refunded exactly once

if funded == false:
  release and refund transitions cannot transfer funds

only buyer can voluntarily release
only seller can voluntarily refund
only arbiter can resolve disputes
\end{lstlisting}

The forbidden states should be visible before code review begins:

{\def\LTcaptype{none} % do not increment counter
\begin{longtable}[]{@{}
  >{\raggedright\arraybackslash}p{(\linewidth - 4\tabcolsep) * \real{0.3333}}
  >{\raggedright\arraybackslash}p{(\linewidth - 4\tabcolsep) * \real{0.3333}}
  >{\raggedright\arraybackslash}p{(\linewidth - 4\tabcolsep) * \real{0.3333}}@{}}
\toprule\noalign{}
\begin{minipage}[b]{\linewidth}\raggedright
Forbidden state
\end{minipage} & \begin{minipage}[b]{\linewidth}\raggedright
Why it matters
\end{minipage} & \begin{minipage}[b]{\linewidth}\raggedright
Required guard
\end{minipage} \\
\midrule\noalign{}
\endhead
\bottomrule\noalign{}
\endlastfoot
\passthrough{\lstinline!released == true!} and
\passthrough{\lstinline!refunded == true!} & Pays both sides from one
escrow & Every payout transition checks both terminal flags \\
\passthrough{\lstinline!funded == false!} and funds transferred out &
Pays from an unfunded escrow & Payouts require successful funding \\
Seller paid more than once & Double spend against escrow balance & Mark
terminal state before unsafe external control flow \\
Buyer refunded more than once & Duplicate refund & Consume
withdrawal/refund state exactly once \\
Arbiter changed by one party during dispute & Authority capture &
Arbiter changes require explicit, delayed, reviewable authority \\
Terminal state reached but transfer failed & Accounting says paid while
value did not move & State change and value movement must have
consistent failure behavior \\
\end{longtable}
}

Now look for broken transitions:

\begin{itemize}
\tightlist
\item
  What if \passthrough{\lstinline!release()!} transfers funds before
  setting \passthrough{\lstinline!released = true!} and the recipient
  can reenter?
\item
  What if \passthrough{\lstinline!arbitrate\_release()!} forgets to
  check \passthrough{\lstinline!refunded == false!}?
\item
  What if \passthrough{\lstinline!fund()!} can be called twice and
  overwrites \passthrough{\lstinline!amount!}?
\item
  What if \passthrough{\lstinline!refund()!} sends funds to
  \passthrough{\lstinline!msg.sender!} instead of
  \passthrough{\lstinline!buyer!}?
\item
  What if the arbiter address can be changed by the seller?
\item
  What if the token used for payment takes a fee on transfer?
\item
  What if the transfer fails but the state is marked released?
\item
  What if neither buyer nor seller cooperates and the arbiter key is
  lost?
\end{itemize}

Notice the range: authorization, accounting, external calls, token
semantics, liveness, and operational risk all appear in a tiny state
machine. The same method scales to larger systems.

\section{Practical Discipline: Do Not Start With Bug Names}

Beginners often ask:

\begin{lstlisting}
Is this vulnerable to reentrancy?
Is this vulnerable to oracle manipulation?
Is this vulnerable to replay?
\end{lstlisting}

Those are useful questions, but they are not the first questions.

Start earlier:

\begin{lstlisting}
What transition changes value?
What authority permits it?
What invariant should it preserve?
What external fact does it rely on?
What happens if the transition occurs twice?
What happens if it occurs before or after another transition?
What happens if it never occurs?
\end{lstlisting}

Then bug classes become labels for specific ways the transition model
can fail.

{\def\LTcaptype{none} % do not increment counter
\begin{longtable}[]{@{}
  >{\raggedright\arraybackslash}p{(\linewidth - 2\tabcolsep) * \real{0.5000}}
  >{\raggedright\arraybackslash}p{(\linewidth - 2\tabcolsep) * \real{0.5000}}@{}}
\toprule\noalign{}
\begin{minipage}[b]{\linewidth}\raggedright
Bug label
\end{minipage} & \begin{minipage}[b]{\linewidth}\raggedright
State-machine interpretation
\end{minipage} \\
\midrule\noalign{}
\endhead
\bottomrule\noalign{}
\endlastfoot
Reentrancy & External control flow allows another transition before
invariants are restored \\
Replay & An input is accepted again after it should have been
consumed \\
Oracle manipulation & A transition depends on an external value whose
production rules do not support the claimed property \\
MEV & Observation, inclusion, or ordering power changes which valid
transition sequence executes \\
Bridge failure & State, proof, finality, domain separation, or trust
assumptions do not line up across machines \\
\end{longtable}
}

This is how the rest of the book will proceed.

\section{Review Questions}

\begin{enumerate}
\def\labelenumi{\arabic{enumi}.}
\tightlist
\item
  What information belongs in ``state'' for a UTXO system and for an
  account-based smart contract system?
\item
  What is the difference between a transaction object and the state
  transition it causes?
\item
  Why can a transaction be valid under explicit rules but unsafe under
  protocol intent?
\item
  Why must replicated blockchain execution be deterministic?
\item
  Give one safety property and one liveness property for a bridge or
  lending protocol.
\item
  Why does ``the chain accepted the transaction'' not prove that the
  application was secure?
\item
  What assumption is being made when a system claims safety below a
  one-third Byzantine threshold?
\end{enumerate}

\section{Applied Exercises}

\subsection{Exercise 1: Model a Token Transfer}

Write a state-machine model for a simple token with:

\begin{itemize}
\tightlist
\item
  \passthrough{\lstinline!balances[address]!}
\item
  \passthrough{\lstinline!total\_supply!}
\item
  \passthrough{\lstinline!transfer(sender, recipient, amount)!}
\item
  \passthrough{\lstinline!mint(recipient, amount)!}
\item
  \passthrough{\lstinline!burn(holder, amount)!}
\end{itemize}

For each transition, write:

\begin{itemize}
\tightlist
\item
  Required authority.
\item
  Validity conditions.
\item
  State updates.
\item
  At least two invariants.
\item
  A transition table.
\end{itemize}

Then answer: which invariants depend on the mint authority being honest
or constrained?

\subsection{Exercise 2: Separate Rule Validity From Intent}

Consider a protocol that allows anyone to deposit collateral and borrow
up to 80 percent of the collateral's oracle value.

The contract checks:

\begin{lstlisting}
collateral_value * 80 / 100 >= borrowed_value
\end{lstlisting}

Assume the oracle price can be manipulated for one block using a
low-liquidity pool.

Answer:

\begin{itemize}
\tightlist
\item
  Is the borrow transaction valid under the contract's explicit rule?
\item
  Which security property is violated?
\item
  What external assumption was false?
\item
  Is this a safety failure, liveness failure, or both?
\end{itemize}

\subsection{Exercise 3: Identify Liveness Risk}

A rollup says users can withdraw through L1 if the sequencer censors
them. The withdrawal path requires a proof generated by an off-chain
service run by the rollup team.

Answer:

\begin{itemize}
\tightlist
\item
  What safety property is the escape hatch supposed to protect?
\item
  What liveness assumption does it introduce?
\item
  What happens if the sequencer is censoring and the proof service is
  down?
\item
  How would you rewrite the security claim more honestly?
\end{itemize}

\subsection{Exercise 4: Broken Transition Taxonomy}

Classify each failure:

\begin{enumerate}
\def\labelenumi{\arabic{enumi}.}
\tightlist
\item
  A bridge accepts the same message twice.
\item
  A validator client accepts a block another client rejects.
\item
  A vault lets the first depositor manipulate share price through direct
  donation.
\item
  A governance proposal executes before the timelock expires.
\item
  A withdrawal queue can become permanently stuck because one recipient
  reverts.
\item
  A frontend asks users to sign a transaction that grants unlimited
  token approval to an attacker.
\end{enumerate}

For each, identify:

\begin{itemize}
\tightlist
\item
  The relevant state.
\item
  The broken transition.
\item
  The missing or false assumption.
\item
  Whether the primary failure is safety, liveness, authorization,
  ordering, or user-side authority.
\end{itemize}

The Part I Integrated Lab will combine this chapter's state-machine
method with assets, authority, threat modeling, and invariants. These
exercises are intentionally smaller: they train the pieces.

\section{Key Takeaways}

\begin{itemize}
\tightlist
\item
  A blockchain is best understood as a replicated state machine
  operating under adversarial inputs.
\item
  State is not only balances. It includes every variable, commitment,
  role, parameter, queue, nonce, proof status, and authority record that
  affects future behavior.
\item
  Transactions are inputs. State transitions are the changes caused by
  applying inputs under rules.
\item
  Validity is layered. A transaction can be syntactically valid, signed,
  included, and executable while still violating protocol intent.
\item
  Replication requires deterministic execution. Honest nodes must
  compute the same result from the same starting state and ordered
  inputs.
\item
  Consensus and fork choice select history. Execution applies transition
  rules to that history.
\item
  Safety means forbidden states or conflicting histories do not happen.
  Liveness means required progress remains possible.
\item
  Fault assumptions are part of the security claim. If they are vague,
  the security claim is vague.
\item
  Most serious application-layer exploits are valid paths through an
  underspecified or incorrectly constrained state machine.
\end{itemize}

\section{Further Reading}

\begin{itemize}
\tightlist
\item
  Leslie Lamport,
  \href{https://lamport.azurewebsites.net/pubs/state-machine.pdf}{``Computation
  and State Machines''}. A concise foundation for thinking about
  computation as state machines and reasoning about invariants.
\item
  Gavin Wood,
  \href{https://ethereum.github.io/yellowpaper/paper.pdf}{``Ethereum: A
  Secure Decentralised Generalised Transaction Ledger''}. The formal
  Ethereum specification, especially the blockchain paradigm and
  state-transition framing.
\item
  ethereum.org,
  \href{https://ethereum.org/developers/docs/transactions/}{``Transactions''}.
  A practical introduction to Ethereum transactions as signed
  state-changing instructions.
\item
  ethereum.org,
  \href{https://ethereum.org/developers/docs/blocks/}{``Blocks''}. A
  practical introduction to blocks as ordered batches of transactions
  linked by hashes.
\item
  Satoshi Nakamoto, \href{https://bitcoin.org/bitcoin.pdf}{``Bitcoin: A
  Peer-to-Peer Electronic Cash System''}. The original Bitcoin paper,
  especially the network, proof-of-work, and double-spend sections.
\item
  Bowen Alpern and Fred B. Schneider, ``Defining Liveness'' and
  ``Recognizing Safety and Liveness.'' These papers are the classical
  formal background for the safety/liveness distinction used throughout
  distributed systems.
\end{itemize}
